% page 584 of the facsimile (image p-583 of the scan)
% block 157.5 x 246.3 mm ; left margin 8.0 mm
\pgc{-12.700mm}{0.000mm}{
\l{\cel{3}576}
\l{}
\l{\sou{teorbo}. (\sou{muziko)}. Instrumento analoga a liuto, plu granda,}
\l{\cel{3}kun du kapi : un por la kordi quin onu fingragas sur la}
\l{\cel{3}mancho, la altra por la grosa kordi quin onu pinchas va-}
\l{\cel{3}kue. - DEFIRS.}
\l{}
\l{\sou{teoremo}. Propoziciono ciencala quan demonstro igas evidenta.}
\l{\cel{3}- DEFIRS.}
\l{}
\l{\sou{teorio}. Ensemblo ek la legi, la reguli atribuita a serio de}
\l{\cel{3}fakti determinita, la principi quin onu propozas kom lia}
\l{\cel{3}existo-kauzo. - DEFIRS.}
\l{}
\l{\sou{teozofio}. I.Sorto di iluminismo religiala. - II. Doktrino}
\l{\cel{3}fondita sur intuico kontinue plu altagrada di la deajo, e}
\l{\cel{3}qua posibligas efikar a la forci superiora. - DEFIRS.}
\l{}
\l{\sou{tepalo}. (\sou{bot.}) Singla ek la parti di la envelopilo florala.}
\l{}
\l{\sou{tepida}. Di qua la temperaturo esas inter kolda e varma.-EFI.}
\l{}
\l{\sou{teplico}. Chambrego vitro-mura, kalorizita, en qua onu shir-}
\l{\cel{3}mas la planti qui domajesus da temperaturo kolda. - R.}
\l{}
\l{\sou{tero.}\cel{2}I. Parto solida di la globo quan ni habitas : (a) sulo}
\l{\cel{3}qua suportas la homo, la animali, la planti, la urbi, e c.;}
\l{\cel{3}(b) sulo qua produktas la vejetanti ; (c) sulo aden qua}
\l{\cel{3}onu sepultas la mortinti. - II. Materio solida ye qua kon-}
\l{\cel{3}sistas nia globo ; substanco tirita de ta materio ed uzata}
\l{\cel{3}diversa-skope. - III. La ter-globo ipsa, kom planeto. -}
\l{\cel{3}EFIS.}
\l{}
\l{\sou{terakoto}. Argilo fasonita e pose bakita. - DEFIRS.}
\l{}
\l{\sou{terapio.(medic.}) Parto di medicino qua havas kom objekto la}
\l{\cel{3}kuraco e risanigo de la morbi. - DEFIRS.}
\l{}
\l{\sou{terario}. Instaluro plu o min granda, en qua onu depozas tero,}
\l{\cel{3}petri, stoni, planti, e c., por la eduko e la entrateno}
\l{\cel{3}di kolektajo ek repteri, batraki, insekti, e c. - EF.}
\l{}
\l{\sou{teraso}. I. Platformo plu o min vasta, facita en gardeno, en}
\l{\cel{3}parko, por regardar la peizajo til horizonto o por prome-}
\l{\cel{3}nar - ek ter-amaso sustenata cirkonde da masonuro. - II.}
\l{\cel{3}Platformo masonura, facita por regardar la cirkondajo o}
\l{\cel{3}por promenar, sur un ek la etaji supra o sur la tekto-}
\l{\cel{3}karpenturo di domo, di edifico. - DEF.}
\l{}
\l{\sou{teratologio}. Parto di la naturo-historio qua deskriptas la}
\l{\cel{3}monstri. -\cel{2}DEF.}
\l{}
\l{terco. I. (\sou{skermo}).Poziciono di la espado kande lu engajesas}
\l{\cel{3}en la lineo extera, kun lua pinto vers-supre, la karpo rot-}
\l{\cel{3}acante de extere ad interne. - II. (muziko) Intervalo inter}
\l{\cel{3}lasekundo e la quarto. - III.(tempo). La parto sisadekima}
\l{\cel{3}di sekundo. - DEFIRS.}
}
